% !TeX program = xelatex
\ifdefined\BookMaster\else
\documentclass[11pt,oneside]{article}
% !TeX program = xelatex
\usepackage[b5paper,left=22mm,right=22mm,top=20mm,bottom=22mm,footskip=11mm]{geometry}
\usepackage{fontspec}
\setmainfont{TeX Gyre Termes}
\setsansfont{TeX Gyre Heros}
\usepackage{amsmath,amssymb,mathtools,amsthm}
\renewcommand{\cdots}{\mathinner{\cdotp\cdotp\cdotp}}
\usepackage{graphicx,adjustbox,array,booktabs,longtable,multirow,enumitem,changepage,needspace}
\usepackage[table]{xcolor}
\usepackage{fancyhdr}
\raggedbottom
\interfootnotelinepenalty=10000
\newcommand{\BookFootnote}[2]{\begingroup\renewcommand{\thefootnote}{\fnsymbol{footnote}}\footnote[#1]{#2}\endgroup}
\newcommand{\BookFootnoteMark}[1]{\begingroup\renewcommand{\thefootnote}{\fnsymbol{footnote}}\footnotemark[#1]\endgroup}
\newcommand{\BookFootnoteText}[2]{\begingroup\renewcommand{\thefootnote}{\fnsymbol{footnote}}\footnotetext[#1]{#2}\endgroup}
\newenvironment{BookFloat}[1]{\par\addvspace{.6\baselineskip}\noindent\def\BookFloatID{#1}\begin{minipage}{\linewidth}\phantomsection\label{bookfloat:#1:start}}{\phantomsection\label{bookfloat:\BookFloatID:end}\end{minipage}\par\addvspace{.6\baselineskip}}

\usepackage[unicode,hidelinks]{hyperref}
\pagestyle{fancy}
\fancyhf{}
\newcommand{\BookFooter}{}
\fancyfoot[L]{\parbox[b]{\dimexpr\textwidth-15mm}{\footnotesize\BookFooter}}
\fancyfoot[R]{\footnotesize\thepage}
\renewcommand{\headrulewidth}{0pt}
\renewcommand{\footrulewidth}{0pt}
\setlength{\parindent}{0pt}
\setlength{\parskip}{0.2em}
\setlength{\emergencystretch}{3em}
\setlist{nosep,leftmargin=1.7em,labelsep=.65em}
\newif\ifBookNarrative
\BookNarrativetrue
\newenvironment{BookBody}{\setlength{\parindent}{0pt}}{}
\newcommand{\BookStatement}[1]{\par\addvspace{.6\baselineskip}\noindent #1\enspace\ignorespaces}
\newcommand{\BookLabel}[1]{\par\addvspace{.6\baselineskip}\needspace{3\baselineskip}\noindent\textbf{#1}\quad\ignorespaces}
\newcommand{\BookHeading}[1]{\par\addvspace{\baselineskip}\needspace{4\baselineskip}\noindent{\large\bfseries #1}\par\medskip}
\newcommand{\BookSubheading}[1]{\par\addvspace{.7\baselineskip}\needspace{3\baselineskip}\noindent{\bfseries #1}\par\smallskip}
\newcommand{\BookTextHeading}[1]{\par\addvspace{.7\baselineskip}\needspace{3\baselineskip}\noindent #1\par\smallskip}
\newcommand{\BookSectionTitle}[2]{\par\addvspace{1.2\baselineskip}\needspace{5\baselineskip}\noindent\colorbox{gray!25}{\sffamily\bfseries\small SECTION #1}\quad{\large\bfseries #2}\par\medskip}
\newcommand{\BookExerciseTitle}[1]{\par\addvspace{\baselineskip}\needspace{4\baselineskip}\noindent{\color{gray}\rule{1.5mm}{1.5mm}}\enspace{\large\bfseries Exercises #1}\par\medskip}
\newcommand{\BookBullet}[1]{\begin{list}{\textbullet}{\setlength{\leftmargin}{1.7em}\setlength{\labelwidth}{1em}\setlength{\labelsep}{.7em}\setlength{\topsep}{0pt}\setlength{\parsep}{0pt}\setlength{\itemsep}{0pt}}\item #1\end{list}}
\newcommand{\BookExercise}[2]{\begin{list}{\textbf{#1}}{\setlength{\leftmargin}{2.2em}\setlength{\labelwidth}{1.5em}\setlength{\labelsep}{.7em}\setlength{\topsep}{.15em}\setlength{\parsep}{0pt}\setlength{\itemsep}{0pt}}\item #2\end{list}}
\newcommand{\BookEnd}[1]{\unskip\nobreak\hfill\ensuremath{#1}}
\newcommand{\BookSource}[2]{\phantomsection\label{sourcepage:#1}} % Source-page anchors, retained through reflow.
\usepackage{multicol}
\usepackage{tikz}
\newcommand{\BookPartEntry}[2]{\par\noindent\makebox[17mm][l]{\bfseries Section #1}#2\par}
\newcommand{\BookPartTitle}[3]{\par\bigskip\noindent\begin{tikzpicture}\node[anchor=north west,inner sep=0pt,outer sep=0pt] (bookpart) at (0,0) {\begin{minipage}{\dimexpr\linewidth-.5pt}\hspace*{26mm}\begin{minipage}[t]{\dimexpr\linewidth-26mm}\begin{adjustbox}{max width=\linewidth}{\fontsize{20}{24}\selectfont\bfseries #2}\end{adjustbox}\par\vspace{6mm}\hrule\vspace{6mm}{\sffamily\footnotesize\setlength{\parskip}{.3em}#3}\vspace{3mm}\end{minipage}\end{minipage}};\path[shade,top color=gray!60,bottom color=white] (0,0) rectangle ([xshift=24mm]bookpart.south west);\node[anchor=north,font=\sffamily\bfseries\small] at (12mm,-2mm) {PART};\node[anchor=north,font=\sffamily\bfseries\Huge] at (12mm,-7mm) {#1};\end{tikzpicture}\par\bigskip}
\newsavebox{\BookMathBox}
\makeatletter
\newcommand{\BookInline}[1]{\leavevmode\hskip0pt\relax\begingroup\in@{\begin}{#1}\ifin@\sbox{\BookMathBox}{\(#1\)}\ifdim\wd\BookMathBox>\dimexpr\linewidth-1em\relax\resizebox{\dimexpr\linewidth-1em\relax}{!}{\usebox{\BookMathBox}}\else\(#1\)\fi\else\in@{\left}{#1}\ifin@\sbox{\BookMathBox}{\(#1\)}\ifdim\wd\BookMathBox>\dimexpr\linewidth-1em\relax\resizebox{\dimexpr\linewidth-1em\relax}{!}{\usebox{\BookMathBox}}\else\(#1\)\fi\else\(#1\)\fi\fi\endgroup\hskip0pt\relax}
\makeatother
\DeclareMathOperator{\Aut}{Aut}
\DeclareMathOperator{\End}{End}
\DeclareMathOperator{\Hom}{Hom}
\DeclareMathOperator{\GF}{GF}
\DeclareMathOperator{\irr}{irr}
\DeclareMathOperator{\lt}{lt}
\DeclareMathOperator{\lp}{lp}
\providecommand{\textcircled}[1]{\textnormal{(#1)}}

\begin{document}
\fi
\renewcommand{\BookFooter}{Appendix: Matrix Algebra}
\pdfbookmark[0]{Appendix: Matrix Algebra}{appendix_matrix_algebra}
\ifdefined\BookMaster\else\setcounter{page}{391}\fi
% Source PDF page 408; original printed page 391
\BookSource{408}{391}
\BookNarrativetrue

\begin{BookBody}

\BookTextHeading{\textbf{Appendix: Matrix Algebra}  }

\noindent We give a brief summary of matrix algebra here. Matrices appear in examples in some chapters of the text and also are involved in several exercises.

A\textbf{ matrix} is a rectangular array of numbers. For example, the array

\begin{center}\begin{adjustbox}{max width=\linewidth}\BookInline{\displaystyle \begin{bmatrix} 2 & -1 & 4 \\ 3 & 1 & 2 \end{bmatrix}}\end{adjustbox}\end{center}

\noindent is a matrix having two rows and three columns. A matrix having \BookInline{m} rows and \BookInline{n} columns is an \BookInline{m\times n} matrix, so Matrix (1) is a \BookInline{2\times 3} matrix. If \BookInline{m= n,\allowbreak } the matrix is\textbf{ square}. Entries in a matrix may be any type of number—integer, rational, real, or complex. We let \BookInline{M_{m\times n}(\mathbb{R})} be the set of all \BookInline{m\times n} matrices with real number entries. If \BookInline{m= n,\allowbreak } the notation is abbreviated to \BookInline{M_{n}(\mathbb{R}).} We can similarly consider \BookInline{M_{n}(\mathbb{Z}),\allowbreak M_{2\times 3}(\mathbb{C}),\allowbreak } etc.

Two matrices having the same number \BookInline{m} of rows and the same number \BookInline{n} of columns can be added in the obvious way: we add entries in corresponding positions.

\textbf{A1 Example}

\noindent In \BookInline{M_{2\times 3}(\mathbb{Z}),\allowbreak } we have

\begin{center}\begin{adjustbox}{max width=\linewidth}\BookInline{\displaystyle \begin{bmatrix} 2 & - 1 & 4 \\ 3 & 1 & 2 \end{bmatrix} + \begin{bmatrix} 1 & 0 & - 3 \\ 2 & - 7 & 1 \end{bmatrix} = \begin{bmatrix} 3 & - 1 & 1 \\ 5 & - 6 & 3 \end{bmatrix}.}\end{adjustbox}\end{center}

We will use uppercase letters to denote matrices. If \BookInline{A,\allowbreak B,\allowbreak } and \BookInline{C} are \BookInline{m\times n} matrices, it is easily seen that \BookInline{A+ B= B+ A} and that \BookInline{A+ (B+ C)= (A+ B)+ C.}

Matrix multiplication\BookInline{,\allowbreak AB,\allowbreak } is defined only if the number of columns of \BookInline{A} is equal to the number of rows of \BookInline{B.} That is, if \BookInline{A} is an \BookInline{m\times n} matrix, then \BookInline{B} must be an \BookInline{n\times s} matrix for some integer \BookInline{s.} We start by defining as follows the product \BookInline{AB} where \BookInline{A} is a \BookInline{1\times n} matrix and \BookInline{B} is an \BookInline{n\times 1} matrix:

\begin{center}\begin{adjustbox}{max width=\linewidth}\BookInline{\displaystyle AB = \begin{bmatrix} a_{1} & a_{2} & \cdots & a_{n} \end{bmatrix} \begin{bmatrix} b_{1} \\ b_{2} \\ \vdots \\ b_{n} \end{bmatrix} = a_{1}b_{1} + a_{2}b_{2} + \cdots + a_{n}b_{n}.}\end{adjustbox}\end{center}

\noindent Note that the result is a number. (We shall not distinguish between a number and the \BookInline{1\times 1} matrix having that number as its sole entry.) You may recognize this product as the\textit{ dot product} of vectors. Matrices having only one\textit{ row} or only one\textit{ column} are\textbf{ row} \textbf{vectors} or\textbf{ column vectors}, respectively.

\end{BookBody}

% Source PDF page 409; original printed page 392
\BookSource{409}{392}
\BookNarrativetrue

\begin{BookBody}

\noindent \textbf{A2 Example} We find that

\begin{center}\begin{adjustbox}{max width=\linewidth}\BookInline{\displaystyle \begin{bmatrix} 3 & - 7 & 2 \end{bmatrix} \begin{bmatrix} 1 \\ 4 \\ 5 \end{bmatrix} = (3)(1) + ( - 7)(4) + (2)(5) = - 15.}\end{adjustbox}\end{center} \BookEnd{\blacktriangle}

Let \BookInline{A} be an \BookInline{m\times n} matrix and let \BookInline{B} be an \BookInline{n\times s} matrix. Note that the number \BookInline{n} of entries in each row of \BookInline{A} is the same as the number \BookInline{n} of entries in each column of \BookInline{B.} The product \BookInline{C= AB} is an \BookInline{m\times s} matrix. The entry in the \BookInline{i}th row and \BookInline{j}th column of \BookInline{AB} is the product of the \BookInline{i}th row of \BookInline{A} times the \BookInline{j}th column of \BookInline{B} as defined by Eq. (2) and illustrated in Example A2.

\noindent \textbf{A3 Example} Compute

\begin{center}\begin{adjustbox}{max width=\linewidth}\BookInline{\displaystyle \boldsymbol { A } \boldsymbol { B } = \begin{bmatrix} 2 & - 1 & 3 \\ 1 & 4 & 6 \end{bmatrix} \begin{bmatrix} 3 & 1 & 2 & 1 \\ 1 & 4 & 1 & - 1 \\ - 1 & 0 & 2 & \quad 1 \end{bmatrix}.}\end{adjustbox}\end{center}

\BookLabel{\textit{Solution}} Note that \BookInline{A} is \BookInline{2\times 3} and \BookInline{B} is \BookInline{3\times 4.} Thus \BookInline{AB} will be \BookInline{2\times 4.} The entry in its second row and third column is

\begin{center}\begin{adjustbox}{max width=\linewidth}\BookInline{\displaystyle (2\mathrm{nd}\;\mathrm{row}\;A)(3\mathrm{rd}\;\mathrm{column}\;B)=[1\quad4\quad6]\begin{bmatrix}2\\1\\2\end{bmatrix}=2+4+12=18.}\end{adjustbox}\end{center}

\noindent Computing all eight entries of \BookInline{AB} in this fashion, we obtain

\begin{center}\begin{adjustbox}{max width=\linewidth}\BookInline{\displaystyle AB = \begin{bmatrix} 2 & -2 & 9 & 6 \\ 1 & 17 & 18 & 3 \end{bmatrix}.}\end{adjustbox}\end{center} \BookEnd{\blacktriangle}

\noindent \textbf{A4 Example} The product

\begin{center}\begin{adjustbox}{max width=\linewidth}\BookInline{\displaystyle \begin{bmatrix} 2 & -1 & 3 \\ 1 & 4 & 6 \end{bmatrix} \begin{bmatrix} 2 & 1 \\ 5 & 4 \end{bmatrix}}\end{adjustbox}\end{center}

\noindent is not defined, since the number of entries in a row of the first matrix is not equal to the number of entries in a column of the second matrix\BookInline{.}\BookEnd{\blacktriangle}

For square matrices of the same size, both addition and multiplication are always defined. Exercise 10 asks us to illustrate the fact that matrix multiplication is not commutative.

That is\BookInline{,\allowbreak AB} need not equal \BookInline{BA} even when both products are defined, as for \BookInline{A,\allowbreak B\in } \BookInline{M_{2}(\mathbb{Z}).} It can be shown that \BookInline{A(BC)= (AB)C} and \BookInline{A(B+ C)= AB+ } \textit{AC} whenever all these expressions are defined.

We let \BookInline{I_{n}} be the \BookInline{n\times n} matrix with entries 1 along the diagonal from the upper-left corner to the lower-right corner, and entries 0 elsewhere. For example,

\begin{center}\begin{adjustbox}{max width=\linewidth}\BookInline{\displaystyle I_{3} = \begin{bmatrix} 1 & 0 & 0 \\ 0 & 1 & 0 \\ 0 & 0 & 1 \end{bmatrix}.}\end{adjustbox}\end{center}

\noindent It is easy to see that if \BookInline{A} is any \BookInline{n\times s} matrix and \BookInline{B} is any \BookInline{r\times n} matrix, then \BookInline{I_{n}A= A} and \textit{BI}\BookInline{n= B.} That is, the matrix \BookInline{I_{n}} acts much as the number 1 does for multiplication when multiplication by \BookInline{I_{n}} is defined.

\end{BookBody}

% Source PDF page 410; original printed page 393
\BookSource{410}{393}
\BookNarrativetrue

\begin{BookBody}

Let \BookInline{A} be an \BookInline{n\times n} matrix and consider a matrix equation of the form\textit{ AX} \BookInline{= B,\allowbreak } where \BookInline{A} and \BookInline{B} are known but \BookInline{X} is unknown. If we can find an \BookInline{n\times n} matrix \BookInline{A^{- 1}} such that \BookInline{A^{- 1}A= AA^{- 1}= I_{n},\allowbreak } then we can conclude that

\begin{center}\begin{adjustbox}{max width=\linewidth}\BookInline{\displaystyle A^{- 1}(AX)= A^{- 1}B,\allowbreak (A^{- 1}A)X= A^{- 1}B,\allowbreak I_{n}X= A^{- 1}B,\allowbreak X= A^{- 1}B,\allowbreak }\end{adjustbox}\end{center}

\noindent and we have found the desired matrix \BookInline{X.} Such a matrix \BookInline{A^{- 1}} acts like the reciprocal of a number\BookInline{:A^{- 1}A= I_{n}} and \BookInline{(1/r)r= 1.} This is the reason for the notation \BookInline{A^{- 1}.}

If \BookInline{A^{- 1}} exists, the square matrix \BookInline{A} is\textbf{ invertible} and \BookInline{A^{- 1}} is the\textbf{ inverse} of \BookInline{A.} If \BookInline{A^{- 1}} does not exist, then \BookInline{A} is said to be\textbf{ singular}. It can be shown that if there exists an \BookInline{n\times n} matrix \BookInline{A^{- 1}} such that \BookInline{A^{- 1}A= I_{n},\allowbreak } then \BookInline{AA^{- 1}= I_{n}} also, and furthermore, there is only one matrix \BookInline{A^{- 1}} having this property.

\noindent \textbf{A5 Example} Let

\begin{center}\begin{adjustbox}{max width=\linewidth}\BookInline{\displaystyle A = \begin{bmatrix} 2 & 9 \\ 1 & 4 \end{bmatrix}.}\end{adjustbox}\end{center}

\noindent We can check that

\begin{center}\begin{adjustbox}{max width=\linewidth}\BookInline{\displaystyle \begin{bmatrix} - 4 & 9 \\ 1 & - 2 \end{bmatrix} \begin{bmatrix} 2 & 9 \\ 1 & 4 \end{bmatrix} = \begin{bmatrix} 2 & 9 \\ 1 & 4 \end{bmatrix} \begin{bmatrix} - 4 & 9 \\ 1 & - 2 \end{bmatrix} = \begin{bmatrix} 1 & 0 \\ 0 & 1 \end{bmatrix}.}\end{adjustbox}\end{center}

\noindent Thus,

\begin{center}\begin{adjustbox}{max width=\linewidth}\BookInline{\displaystyle A^{-1} = \begin{bmatrix} -4 & 9 \\ 1 & -2 \end{bmatrix}.}\end{adjustbox}\end{center} \BookEnd{\blacktriangle}

We leave the problems of determining the existence of \BookInline{A^{- 1}} and its computation to a course in linear algebra.

Associated with each square \BookInline{n\times n} matrix \BookInline{A} is a number called the\textit{ determinant} of \BookInline{A} and denoted by det\BookInline{(A).} This number can be computed as sums and differences of certain products of the numbers that appear in the matrix \BookInline{A.} For example, the determinant of the \BookInline{2\times 2} matrix \BookInline{\left[ \begin{matrix} { a \quad b } \\ { c \quad d } \end{matrix} \right]} is \BookInline{ad- bc.} Note that an \BookInline{n\times 1} matrix with real number entries can be viewed as giving coordinates of a point in \BookInline{n-}dimensional Euclidean space \BookInline{\mathbb{R}^{n}.} Multiplication of such a single column matrix on the left by a real \BookInline{n\times n} matrix \BookInline{A} produces another such single column matrix corresponding to another point in \BookInline{\mathbb{R}^{n}.} This multiplication on the left by \BookInline{A} thus gives a map of \BookInline{\mathbb{R}^{n}} into itself. It can be shown that a piece of \BookInline{\mathbb{R}^{n}} of volume \BookInline{V} is mapped by this multiplication by \BookInline{A} into a piece of volume \BookInline{|}det\BookInline{(A)|\cdot V.} This is one of the reasons that determinants are important.

The following properties of determinants for \BookInline{n\times n} matrices \BookInline{A} and \BookInline{B} are of interest in this text:

\BookExercise{1.}{det\BookInline{(I_{n})= 1}}

\BookExercise{2.}{det\BookInline{(AB)= } det\BookInline{(A)} det\BookInline{(B)}}

\BookExercise{3.}{det\BookInline{(A)\neq 0} if and only if \BookInline{A} is an invertible matrix}

\BookExercise{4.}{If \BookInline{B} is obtained from \BookInline{A} by interchanging two rows (or two columns) of \BookInline{A,\allowbreak } then det\BookInline{(B)= - } det\BookInline{(A)}}

\BookExercise{5.}{If every entry of \BookInline{A} is zero above the\textit{ main diagonal} from the upper left corner to the lower right corner, then det \BookInline{(A)} is the product of the entries on this diagonal. The same is true if all entries below the main diagonal are zero.}

\end{BookBody}

% Source PDF page 411; original printed page 394
\BookSource{411}{394}
\BookNarrativefalse

\begin{BookBody} \BookEnd{\blacksquare}

\BookTextHeading{EXERCISES A  }

\noindent In Exercises 1 through 9, compute the given arithmetic matrix expression, if it is defined.

\begin{center}\begin{adjustbox}{max width=\linewidth}\BookInline{\displaystyle \begin{bmatrix} - 2 & 4 \\ 1 & 5 \end{bmatrix} + \begin{bmatrix} 4 & - 3 \\ 1 & 2 \end{bmatrix}}\end{adjustbox}\end{center}

\begin{center}\begin{adjustbox}{max width=\linewidth}\BookInline{\displaystyle i.\ \begin{bmatrix} 1 + i & - 2 & 3 - i \\ 4 & i & 2 - i \end{bmatrix} + \begin{bmatrix} 3 & i - 1 & - 2 + i \\ 3 - i & 1 + i & 0 \end{bmatrix}}\end{adjustbox}\end{center}

\noindent \textbf{2}

\begin{center}\begin{adjustbox}{max width=\linewidth}\BookInline{\displaystyle \begin{bmatrix} i & - 1 \\ 4 & 1 \\ 3 & - 2i \end{bmatrix} - \begin{bmatrix} 3 - i & 4i \\ 2 & 1 + i \\ 3 & - i \end{bmatrix}}\end{adjustbox}\end{center}

\begin{center}\begin{adjustbox}{max width=\linewidth}\BookInline{\displaystyle 4.\ \begin{bmatrix} 1 & -1 \\ 3 & \quad 1 \end{bmatrix} \begin{bmatrix} \quad 2 & 4 \\ -1 & 3 \end{bmatrix}}\end{adjustbox}\end{center}

\begin{center}\begin{adjustbox}{max width=\linewidth}\BookInline{\displaystyle 5.\ \begin{bmatrix} 3 & 1 \\ -4 & 2 \end{bmatrix} \begin{bmatrix} 1 & 5 & -3 \\ 2 & 1 & 6 \end{bmatrix}}\end{adjustbox}\end{center}

\begin{center}\begin{adjustbox}{max width=\linewidth}\BookInline{\displaystyle \begin{bmatrix} 4 & - 1 \\ 3 & 2 \end{bmatrix} \begin{bmatrix} 1 & 0 \\ - 1 & 7 \\ 3 & 1 \end{bmatrix}}\end{adjustbox}\end{center}

\begin{center}\begin{adjustbox}{max width=\linewidth}\BookInline{\displaystyle 7.\ \begin{bmatrix} \boldsymbol{i} & 1 \\ -2 & 1 \end{bmatrix} \begin{bmatrix} 3\boldsymbol{i} & 1 \\ 4 & -2\boldsymbol{i} \end{bmatrix}}\end{adjustbox}\end{center}

\begin{center}\begin{adjustbox}{max width=\linewidth}\BookInline{\displaystyle 8.\ \begin{bmatrix} 1 & -1 \\ 1 & 0 \end{bmatrix}^{4}}\end{adjustbox}\end{center}

\begin{center}\begin{adjustbox}{max width=\linewidth}\BookInline{\displaystyle \begin{bmatrix} 1 & - i \\ i & 1 \end{bmatrix}^{4}}\end{adjustbox}\end{center}

\BookExercise{10.}{Give an example in \BookInline{M_{2}(\mathbb{Z})} showing that matrix multiplication is not commutative.}

\BookInline{\Big[\Big]^{- 1}}

\BookExercise{11.}{Find \BookInline{\begin{bmatrix} 0 & 1 \\ -1 & 0 \end{bmatrix}^{-1}}, by experimentation if necessary.}

\BookInline{\begin{bmatrix} 2 & 0 & 0 \\ 0 & 4 & 0 \\ 0 & 0 & - 1 \end{bmatrix}^{- 1}}

\BookExercise{12.}{Find \BookInline{\begin{bmatrix} 2 & 0 & 0 \\ 0 & 4 & 0 \\ 0 & 0 & - 1 \end{bmatrix}^{- 1}}, by experimentation if necessary.}

\BookInline{\mathrm{A} = \begin{bmatrix} 3 & 0 & 0 \\ 10 & -2 & 0 \\ 4 & 17 & 8 \end{bmatrix}}

\BookExercise{13.}{If \BookInline{\mathrm{A} = \begin{bmatrix} 3 & 0 & 0 \\ 10 & -2 & 0 \\ 4 & 17 & 8 \end{bmatrix}}, find det (A).}

\BookExercise{14.}{Prove that if \BookInline{A,\allowbreak B\in M_{n}(\mathbb{C})} are invertible, then \BookInline{AB} and \BookInline{BA} are invertible also.}

\end{BookBody}
\ifdefined\BookMaster\let\BookFinish\relax\else\def\BookFinish{\end{document}}\fi\BookFinish
